% Tabla A3-T1. Generada automáticamente; no editar a mano.
% Requiere \usepackage{booktabs} (y graphicx si la tabla es ancha).
\begin{table}[htbp]
    \centering
    \caption{Análisis III: construcción de la muestra de sistemas multiplanetarios.}
    \label{tab:A3-T1}
    \small
    \begin{tabular}{p{0.40\textwidth}rr}
    \toprule
    Paso & $N_{\rm pl}$ & $N_\star$ \\
    \midrule
    Catálogo limpio & 6\,267 & 4\,711 \\
    Con radio, período, {[Fe/H]} y masa estelar & 5\,578 & 4\,105 \\
    Radio planetario entre 1 y 4 $R_\oplus$ & 3\,371 & 2\,419 \\
    Período orbital menor que 100 días & 3\,238 & 2\,340 \\
    En sistemas con al menos 2 planetas que cumplen los criterios & 1\,513 & 615 \\
    \bottomrule
    \end{tabular}
    \par\smallskip
    {\footnotesize\textit{Nota:} Las razones de período se calculan solo entre los planetas que cumplen los criterios, ordenados por período.}
\end{table}
